\documentclass[11pt, a4paper, twoside, titlepage]{scrbook}

% ==========================================
% KOMA-SCRIPT FONT SETTINGS
% ==========================================
% Force serif font (classic look) for all headings, labels, and headers/footers
\setkomafont{disposition}{\normalfont\bfseries}
\setkomafont{descriptionlabel}{\normalfont\bfseries}
\setkomafont{pageheadfoot}{\normalfont\small}

% ==========================================
% PACKAGES
% ==========================================
\usepackage[utf8]{inputenc}
\usepackage[T1]{fontenc}
\usepackage[english]{babel}
\usepackage{multicol}

% Matematica e Informatica Teorica
\usepackage{amsmath, amssymb, amsthm, mathtools}
\usepackage{stmaryrd} % Simboli aggiuntivi (es. parentesi per semantica)

% Layout e Geometria
\usepackage{geometry}
% Usiamo inner/outer per il fronte-retro e un offset per la rilegatura
\geometry{a4paper, inner=2.5cm, outer=2.5cm, bindingoffset=0.5cm, top=2.5cm, bottom=2.5cm}
\usepackage{float} % Aggiunto per supportare [H] nelle figure
\usepackage{ccicons} % Added for CC icons

% Grafica e Colori
\usepackage{graphicx}
\usepackage{xcolor}
\usepackage{tikz} % Per grafi, automi, architetture
\usetikzlibrary{automata, positioning, arrows, arrows.meta}

% ==========================================
% CUSTOMIZATION & BRANDING (UNITN STYLE)
% ==========================================
% Colore rosso UniTn (approssimato per uso digitale/stampa)
\definecolor{unitnred}{RGB}{178, 35, 52}
\definecolor{codebg}{RGB}{245, 245, 245}

% ==========================================
% HEADER & FOOTER (KOMA-Script Native)
% ==========================================
\usepackage[headsepline=0.4pt, footsepline=0.4pt, plainfootsepline]{scrlayer-scrpage}
\usepackage{fontawesome5}
\clearpairofpagestyles

% Attiva i marker dinamici: Capitolo a sinistra (pari), Sezione a destra (dispari)
\automark[section]{chapter}

% Intestazioni (ohead = outer, ihead = inner)
\ohead{\headmark} % Nome del capitolo/sezione corrente
\ihead{\textbf{Security Testing}} % Nome del corso fisso all'interno

% Piè di pagina
% L'argomento opzionale [...] definisce lo stile per le pagine dei capitoli (plain)
\ifoot[A.Y. 2026/2027 -- \textit{Notes by Luca Facchini}]{A.Y. 2026/2027 -- \textit{Notes by Luca Facchini}}
\ofoot[Page \pagemark]{Page \pagemark}

% Customizzazione linee (entrambe in rosso UniTn per simmetria)
\addtokomafont{headsepline}{\color{unitnred}}
\addtokomafont{footsepline}{\color{unitnred}}
\setlength{\headheight}{24pt}

% --- MODERN CHAPTER HEADINGS ---
\makeatletter
\renewcommand*{\chapterlinesformat}[3]{%
  \Ifstr{#1}{chapter}{%
    \parbox[b]{\textwidth}{%
      \raggedright
      \Ifstr{#2}{}{}{% If there is a chapter number
        {\fontsize{70}{75}\selectfont\color{unitnred!20}#2}\\[-2ex]
      }%
      \Huge\bfseries\sffamily\color{black}#3\par\vspace{1cm}
    }%
  }{% For unnumbered chapters (like Preface/Appendix)
    \@hangfrom{#2}{#3}%
  }%
}
\makeatother

% ==========================================
% BOXES, CODE & ALGORITHMS
% ==========================================
\usepackage{tcolorbox}
\tcbuselibrary{theorems, breakable, skins}

\usepackage{listings}
\usepackage{algorithm}
\usepackage{algpseudocode}

% Stile del Codice
\lstset{
    backgroundcolor=\color{codebg},
    basicstyle=\ttfamily\footnotesize,
    keywordstyle=\color{unitnred}\bfseries,
    commentstyle=\color{green!60!black}\itshape,
    stringstyle=\color{blue},
    breaklines=true,
    frame=single,
    rulecolor=\color{lightgray},
    numbers=left,
    numberstyle=\tiny\color{gray},
    showstringspaces=false
}

% ==========================================
% LINKS, REFERENCING & GLOSSARIES
% ==========================================
\usepackage{hyperref}
\hypersetup{
    colorlinks=true,
    linkcolor=unitnred,
    filecolor=magenta,      
    urlcolor=unitnred,
    citecolor=blue,
    pdftitle={Lecture Notes - Security Testing},
    pdfauthor={Luca Facchini},
}

\usepackage[acronym,toc]{glossaries-extra}
\makeglossaries

\makeatletter
% 1. Silently and GLOBALLY count EACH time ANY link-generating acronym command is used
\appto\glslinkpostsetkeys{%
    \ifcsundef{acrocount@\glslabel}{%
        \csgdef{acrocount@\glslabel}{1}%
    }{%
        % \csxdef forces the math to be saved globally, surviving any local groups
        \csxdef{acrocount@\glslabel}{\the\numexpr\csuse{acrocount@\glslabel} + 1\relax}%
    }%
}

% 2. Use the native glossaries-extra hook to append the count in the printed glossary
\renewcommand{\glsxtrpostdescacronym}{%
    \ifcsdef{acrocount@\glscurrententrylabel}{%
        \ifnum\csuse{acrocount@\glscurrententrylabel}=1\relax
            \space \textit{[1]}%
        \else
            \space \textit{[\csuse{acrocount@\glscurrententrylabel}]}%
        \fi
    }{}%
}
\makeatother

\usepackage{cleveref}

\usepackage{microtype}
% --- PREMIUM FONTS ---
\usepackage{newpxtext, newpxmath} % Palatino font for body and math
\usepackage[scaled=0.95]{inconsolata} % Sleek, modern monospace font for code
\usepackage[scaled=0.92]{helvet} % Clean sans-serif for headings
% Set all section headings to use the sans-serif font
\setkomafont{disposition}{\sffamily\bfseries}

% ==========================================
% ACRONYMS DEFINITIONS
% ==========================================
\newacronym{cia}{CIA}{Confidentiality, Integrity, and Availability}
\newacronym{cwe}{CWE}{Common Weakness Enumeration}
\newacronym{cve}{CVE}{Common Vulnerabilities and Exposures}
\newacronym{cvss}{CVSS}{Common Vulnerability Scoring System}
\newacronym{nvd}{NVD}{National Vulnerability Database}
\newacronym{euvd}{EUVD}{European Union Vulnerability Database}
\newacronym{enisa}{ENISA}{European Union Agency for Cybersecurity}
\newacronym{os}{OS}{Operating System}
\newacronym{fs}{FS}{File System}
\newacronym{sto}{STO}{Security Through Obscurity}
\newacronym{cert}{CERT}{Computer Emergency Response Team}
\newacronym{sdlc}{SDLC}{System Development Life Cycle}
\newacronym{ssdlc}{sSDLC}{Secure System Development Life Cycle}
\newacronym{qa}{QA}{Quality Assurance}
\newacronym{srs}{SRS}{Software Requirements Specification}
\newacronym{dds}{DDS}{Detailed Design Specification}
\newacronym{uaf}{UAF}{Use-After-Free}
\newacronym{capec}{CAPEC}{Common Attack Pattern Enumeration and Classification}
\newacronym{sast}{SAST}{Static Application Security Testing}
\newacronym{hmac}{HMAC}{Hash-based Message Authentication Code}
\newacronym{owasp}{OWASP}{Open Worldwide Application Security Project}
\newacronym{xss}{XSS}{Cross-Site Scripting}
\newacronym{sql}{SQL}{Structured Query Language}
\newacronym{sqli}{SQLi}{\gls{sql} Injection}
\newacronym{api}{API}{Application Programming Interface}
\newacronym{csrf}{CSRF}{Cross-Site Request Forgery}
\newacronym{ssrf}{SSRF}{Server-Side Request Forgery}
\newacronym{sbom}{SBOM}{Software Bill of Materials}
\newacronym{gdpr}{GDPR}{General Data Protection Regulation}
\newacronym{pcidss}{PCI DSS}{Payment Card Industry Data Security Standard}
\newacronym{orm}{ORM}{Object-Relational Mapping}
\newacronym{ux}{UX}{User Experience}

% ==========================================
% TCOLORBOX THEOREMS & DEFINITIONS
% ==========================================
\newtcbtheorem[number within=chapter,
  crefname={definition}{definitions},
  Crefname={Definition}{Definitions},
  list inside=deflist
]{definition}{Definition}{%
  colback=unitnred!5,colframe=unitnred!80!black,fonttitle=\bfseries, breakable%
}{def}

\newtcbtheorem[number within=chapter,
  crefname={theorem}{theorems},
  Crefname={Theorem}{Theorems},
  list inside=thmlist
]{theorem}{Theorem}{%
  colback=blue!5,colframe=blue!80!black,fonttitle=\bfseries, breakable%
}{thm}

\newtcbtheorem[number within=chapter,
  crefname={lemma}{lemmas},
  Crefname={Lemma}{Lemmas},
  list inside=lemlist
]{lemma}{Lemma}{%
  colback=cyan!5,colframe=cyan!80!black,fonttitle=\bfseries, breakable%
}{lem}

\newtcbtheorem[number within=chapter,
  crefname={observation}{observations},
  Crefname={Observation}{Observations},
  list inside=obslist
]{observation}{Observation}{%
  colback=orange!5,colframe=orange!80!black,fonttitle=\bfseries, breakable%
}{obs}

\newtcbtheorem[number within=chapter,
  crefname={concept}{concepts},
  Crefname={Concept}{Concepts},
  list inside=cptlist
]{concept}{Concept}{%
  colback=green!5!white,colframe=green!60!black,fonttitle=\bfseries, breakable%
}{cpt}

\newcommand{\lecturestop}[1]{%
    \par\vspace{1em}
    \begin{tcolorbox}[
        colback=gray!10,
        colframe=gray!60,
        coltext=gray!60!black,
        fontupper=\centering\bfseries\sffamily,
        boxrule=0.5pt,
        arc=2mm,
        left=4mm, right=4mm, top=2mm, bottom=2mm
    ]
        Lecture stopped here on #1 — Remaining slides to be covered in the next session.
    \end{tcolorbox}
    \par\vspace{1em}
}

% ==========================================
% TITLE PAGE INFO
% ==========================================
\begin{document}
\frontmatter

% --- TITLE PAGE ---
\begin{titlepage}
    % Full-page background graphics using Ti*k*Z
    \begin{tikzpicture}[remember picture, overlay]
        % 1. Bold UniTn red vertical bar on the left edge
        \fill[unitnred] (current page.south west) rectangle ([xshift=1.5cm]current page.north west);
        % 2. Subtle gray geometric accent at the bottom right
        \fill[gray!10] (current page.south east) -- ([yshift=6cm]current page.south east) -- ([xshift=-8cm]current page.south east) -- cycle;
        
        % 3. Giant faint security watermark (Requires \usepackage{fontawesome5} in preamble)
        \node[text=gray!10, font=\fontsize{250}{250}\selectfont, anchor=center] at (current page.center) {\faShield*};
    \end{tikzpicture}
    
    % Push content down and to the right to clear the red bar
    \vspace*{2cm}
    \noindent\hspace{1cm}%
    \begin{minipage}{\dimexpr\textwidth-1cm\relax}
        
        % University Logo
        \includegraphics[width=4.5cm]{UniTn.jpg} 
        
        \vspace{4cm}
        
        % Main Title Area
        {\Huge \bfseries \sffamily \color{unitnred} Security Testing\par}
        \vspace{0.5cm}
        {\LARGE \itshape Lecture Notes\par}
        \vspace{0.5cm}
        {\Large A.Y. 2026/2027\par}
        
        \vfill
        
        % Author Info Block
        \begin{flushleft}
            {\Large \textbf{Luca Facchini}}\\[0.2cm]
            \textit{Master's Degree in Computer Science}\\
            University of Trento\\[0.3cm]
            \href{mailto:luca.facchini-1@studenti.unitn.it}{\texttt{luca.facchini-1@studenti.unitn.it}}\\
            \href{mailto:luca@fc-software.it}{\texttt{luca@fc-software.it}}
        \end{flushleft}
        
        \vspace{1.5cm}
        
        % Footer line and date
        \textcolor{gray}{\rule{\linewidth}{0.8pt}}
        \vspace{0.2cm}\\
        {\small \textbf{Last updated:} \today \hfill \textit{University of Trento}}
        
    \end{minipage}
\end{titlepage}
\clearpage
% --- COPYRIGHT / LEGAL PAGE (Back of Title Page) ---
\newpage
\thispagestyle{empty}
\vspace*{\fill} % Pushes the text to the bottom of the page

\noindent \textbf{About this Document}\\
These notes cover the topics of the \textbf{Security Testing} course, held by Prof. Alessandro Marchetto and TA Dott. Michele Bof during the Master's Degree in Computer Science at the University of Trento. The content herein is derived from class lectures, course recordings, and the professors' official slides and notes.

\vspace{1em}
\noindent \textbf{License \& Copyright}\\
\textcopyright\ 2026 Luca Facchini. My personal contributions, transcriptions, and formatting within this document are distributed under the Creative Commons Attribution-NonCommercial 4.0 International license (\textbf{CC BY-NC 4.0} \ccbynceu). The full text of the license is available in the \texttt{LICENSE} file at the root of this project. Please note that any original intellectual property, figures, or direct excerpts sourced from the professors remain strictly under their respective original copyrights and licenses.

\vspace{1em}
\noindent \textbf{Disclaimer}\\
This material is provided ``AS IS'', without warranty of any kind, express or implied, including but not limited to the warranties of accuracy, completeness, or fitness for a particular purpose. These are independently produced student notes and do not constitute an official university publication. 

\vspace{1em}
\noindent \textit{If you find any typos or errors, please contact me at the email on the title page.}
\vspace{2cm}
\cleardoublepage

% --- TABLE OF CONTENTS ---
% Color the Chapter entries in the TOC
\addtokomafont{chapterentry}{\color{unitnred}}
% Color the Chapter page numbers in the TOC
\addtokomafont{chapterentrypagenumber}{\color{unitnred}}
\tableofcontents
\cleardoublepage
\mainmatter

% --- PREFACE / SCOPE ---
\addchap{Preface: Scope and Purpose}

\noindent \textbf{The Motivation}\\
These notes were initially born out of a personal need: the desire to organize, synthesize, and rigorously typeset the vast amount of information presented during the lectures. Writing and structuring this material in \LaTeX\ served as my primary method of active recall and study. However, recognizing how challenging it can be to study from scattered slides, recordings, and fragmented notes, I decided to open-source this project. The goal is to provide current and future students of the Master's Degree in Computer Science with a unified, structured, and accessible study companion.

\vspace{1em}
\noindent \textbf{The Scope}\\
The scope of this document is to cover the core theoretical concepts, definitions, vulnerabilities, and software life-cycle processes discussed throughout the course. It is designed to act as a comprehensive bridge between the professors' official slides and the verbal explanations provided during class. 

\vspace{1em}
\noindent \textbf{How to Use This Material}\\
While I have made every effort to ensure the technical rigor of this document, it is important to state what these notes are \textit{not}. They are \textbf{not a replacement} for the official recommended textbooks, nor are they a substitute for attending lectures and engaging with the professors. They are a student-made study guide. You should use them to clarify complex topics, review for exams, and quickly reference formal definitions.

\vspace{1em}
\noindent \textbf{Collaborative Spirit}\\
Because this is an actively maintained, open-source project, it is subject to continuous improvement. If you are studying from this document and spot a typographical error, a technical inaccuracy, or a missing concept, please do not hesitate to reach out via the contact information on the title page. Collaborative corrections are highly encouraged and deeply appreciated.

\vspace{2em}
\noindent \textit{Happy studying,}\\
\textbf{Luca Facchini}\\
BSc Graduate in Computer Science

\cleardoublepage

% ==========================================
% LECTURES BEGIN
% ==========================================


\chapter{Introduction to the Course [10/09/2026]}
The objective of the course is to provide principles, techniques, and tools for Security Testing.

The project, as part of the exam, is written in Java, but the course also explores other programming languages like C/C++, Python, PHP, and SQL. This is because different programming languages suffer from different vulnerabilities.

\section{Contents of the Course}
The main topics of the course are described below:

\paragraph{Introduction to Security Testing} An overview of the security testing life-cycle, vulnerabilities, bugs, and secure coding.
\paragraph{Context} A context regarding vulnerability analysis, rates, and repositories will be provided.
\paragraph{Vulnerabilities} A detailed presentation of SQL Injection, XSS, String formatting issues, and more. 
\paragraph{Vulnerability Detection} Both static and dynamic analysis will be discussed, and an overview of AI-based vulnerability detection will be provided.
\paragraph{Tools} For these analyses, tools need to be provided; they will be analyzed and discussed.
\paragraph{[OPTIONAL] Advanced Topics} 

\section{Exam}
The exam is composed of two parts:
\begin{itemize}
    \item A written test with exercises as well as open and closed-form questions (maximum two (2) hours) - About 60\% of the grade.
    \item A written project report \textbf{to be submitted before the written exam} - About 40\% of the grade.
\end{itemize}
It is permitted to attend both in one session or in distinct sessions with no expiration date. Both parts need to be at least 18/30; grades are expressed out of 30.

A simulation of the written part will be performed. Some laboratories will be devoted to working on the exam project.

\subsection{Project Rules}
\begin{itemize}
    \item Groups can be composed of one or two students.
    \item The project can be submitted a maximum of two non-consecutive times.
    \item If the project is graded as a fail two times, a different test has to be arranged.
    \item A passing project can be used for a maximum of two written tests.
    \item In case of a passing final grade, the project cannot be resubmitted.
    \item In case of resubmission, a clear difference/delta between the first and the second submission must exist; otherwise, the final score will be reduced by 1.
\end{itemize}

\textbf{TOTAL OR PARTIAL PLAGIARISM IS FORBIDDEN}

\noindent This applies to both the project and the written test.
\chapter{Concepts, Notions, and Context [10-14/09/2026]}
We will start by defining Software and Security.

\begin{definition}{Software}{software}
    Software is composed of a set of programs and data which provide functionality. Often, it can be part of a cyber-physical system.
\end{definition}

\begin{definition}{Security}{security}
    Understanding and identifying software-induced security risks and determining how to manage them.
\end{definition}

With these definitions, we can conclude that software plays a crucial role in providing security, but it is also a significant \textbf{source of security issues} and problems. Meanwhile, the relevance of security has been increasing but still not enough. Indeed, programming courses do not pay much attention to it, and many developers have limited training in software security or no training at all. This is exacerbated by the fact that the primary goal is usually just to develop the software.

\paragraph{Why is software a source of security vulnerabilities/issues?}
Most developers produce \textbf{insecure code}. Some relevant factors are:
\begin{itemize}
    \item Limited training.
    \item Textbooks do not focus on security.
    \item Limited security audits.
    \item Developers focus on feature implementation and not security.
    \item Security is complex and expensive, requiring time and resources.
    \item Legacy software is often unsafe.
\end{itemize}

Security is frequently considered a \textbf{secondary factor}; often it is a \textbf{trade-off (or conflict)} between: \textbf{Security}, \textbf{Functionality}, and Convenience/Opportunities.

\begin{definition}{Secure Software}{secure_software}
    A software system is \textbf{secure if it satisfies specified security objectives}.

    E.g., \gls{cia} requirements for the system's data and functionality.
\end{definition}

In general, the \gls{cia} triad is not enough; there are other properties like Authentication, Authorization, Privacy, Anonymity, Non-repudiation, Reliability, Audit, etc. \\
Security is about \textbf{preventing risks and imposing countermeasures to reduce threats} to relevant assets to acceptable levels.

\begin{concept}{Security Policy}{security-policy}
    A security policy specifies what security requirements and goals the countermeasures are intended to achieve.
\end{concept}

\begin{concept}{Security Mechanisms}{security_mechanisms}
    Security mechanisms and controls are applied to enforce the security policy as defined in \cref{cpt:security-policy}.
\end{concept}

\begin{concept}{Security by Design}{sec_by_design}
    Security by Design refers to when security is addressed during the design, implementation, and testing phases (i.e., each phase) of software development.
\end{concept}

A system is only as secure as its weakest element. A typical system is composed of several actors:
\begin{description}
    \item[System] Provides data and services.
    \item[Honest user] Uses the system as intended.
    \item[Dishonest attacker] Attempts to find issues using unusual interactions and features to produce harmful effects. Such as disrupting an honest user's access (Integrity and Availability), learning information intended for honest users only (Confidentiality), or intercepting communications (Confidentiality).
    \item[Defender] Protects against possible attacks (known and unknown).
\end{description}

It is important to distinguish between safety and security:
\begin{concept}{Safety vs Security}{safety_security}
    Safety is about protecting from accidental risks, while security is about mitigating \textbf{risks of dangers caused by intentional, malicious actions}.
\end{concept}

\begin{concept}{Security Software Failure}{sec_software_failure}
    A Security software failure is a scenario in which the software does not achieve its security objectives/requirements. Software issues that can cause such failures exist.
\end{concept}

Most software has no clear and specific absolute objectives; rather, these objectives are often specific and case-dependent. Usually, they are traded off against non-functional requirements. For this reason, software design and implementation issues (e.g., bugs) can lead to a \textbf{substantial disruption of the software} and its behavior.

Therefore, \textbf{Security} is a \textbf{system property, not a feature}, and it is a core part of reliability.

\lecturestop{10/09/2026}

\paragraph{Design and Implementation Issues}
There are multiple kinds of implementation issues, such as: \textbf{Bugs in software and environments} (like in architecture, infrastructure, components, etc.), \textbf{inappropriate features} in the infrastructure or reused software (like libraries), and \textbf{inappropriate use} of features provided by the infrastructure or reused software. Possible causes of these are mainly: Unclear or missing requirements, complexity of features, time pressure, characteristics of the programming language, and developer/manager errors.

\begin{definition}{Bug}{bug}
    The system isn't behaving as it's designed to behave.
\end{definition}

\begin{definition}{Vulnerability}{vulnerability}
    A chain of \textbf{weaknesses} linked by causality; it starts with a bug (\cref{def:bug}) and ends with an error which, if exploited, leads to a security failure (\cref{cpt:sec_software_failure}).
\end{definition}

Many security issues are related to vulnerabilities due to bugs, defects, or design flaws in a targeted software system.

\paragraph{Common Weakness Enumeration (\gls{cwe})}
\gls{cwe} provides a publicly available list of the most common software (and hardware) weaknesses that can contribute to security vulnerabilities. Each \gls{cwe} is described by several fields. Why is this done? If the needed function is in a package affected by a known \gls{cwe}, the weakness can be isolated and mitigated during the design of the application before deployment. Also, fixing a weakness before a vulnerability is introduced is much more cost-effective. 

\paragraph{Common Vulnerabilities and Exposures (\gls{cve})}
\gls{cve} is a publicly available list of entries that describes vulnerabilities in widely-used software components. Each \gls{cve} is described by several fields.

\paragraph{\gls{nvd} and \gls{euvd}} 
The \gls{nvd} is a publicly available database with \gls{cve} entries. This database contains a large number of vulnerabilities and provides methods to search and explore the stored \glspl{cve}. The equivalent in Europe is the \gls{euvd}; this is part of the NIS2 directive. The EU mandated \gls{enisa} to develop and maintain the \textbf{EU}ropean \textbf{V}ulnerability \textbf{D}atabase. \gls{enisa} cooperates with MITRE and other non-EU operators working with the \gls{cve} system, putting together various sources to enrich it.

\section{Bug Exploitation}
While a typical bug might just lead to a software crash, an attacker (who is \textbf{not} a normal user) will actively attempt to find bugs/defects through unusual/unexpected interactions and features. They will work to exploit them to do much worse.

\begin{definition}{Exploitation}{exploitation}
    An activity in which an attacker uses a bug/defect (vulnerability, see \cref{def:vulnerability}) to achieve their goals.
\end{definition}

\begin{definition}{Exploit}{exploit}
    A cleverly crafted input, or a series of (unusually unintuitive) interactions with the system, that triggers a bug or a flaw to benefit the attacker.
\end{definition}

As per the previous definitions, an attacker might need several steps to achieve their attack goals; these steps are called an attack path.

\begin{definition}{Attack Path}{attack-path}
    An attack path is a series of exploits (\cref{def:exploit}) or atomic attacks.
\end{definition}

Usually, an exploit can be very complex, but some patterns can be identified. CAPEC is trying to enumerate and classify the most common attacks. 
Also, exploit databases exist that collect \gls{cvss}-compliant exploits and vulnerable software (e.g., ATT\&CK).

\section{Achieving Security}
Security in computer software systems is hard and complicated due to great complexity and dependency on the \gls{os}, \gls{fs}, Network, and physical access. Also, system/software security is \textbf{difficult to measure}; there are no perfect metrics and not enough powerful security testing tools and procedures. From a business perspective, there are deadlines to meet, clients do not demand security (it's often not considered a \textbf{user requirement}), and finally, few to no developers have a culture and knowledge regarding security.

\paragraph{Software Size} 
From a technical POV, the use of software is growing, and along with it, its complexity and size. Also, the reuse and maintenance of pieces of code are growing. Both lead to more bugs and flaws since complex software is difficult to debug and audit for bugs. While it's more probable to find a vulnerability, it is also more difficult and time-consuming for an attacker to find a way in.

\paragraph{Roles and Responsibilities} 
Who is responsible for software security? In our view, Developers and Pen-testers.

\paragraph{How much security?} 
Since absolute security is \textbf{unachievable}, a \textbf{trade-off} must be made:
\begin{center}
    More Security \(\Rightarrow\) More Cost \(\Rightarrow\) Less convenience 
\end{center}
Indeed, security measures (see \cref{cpt:security-policy}) should be as invisible as possible to the end-user, and an appropriate level relevant to the need should be chosen.

\begin{figure}[!htbp]
    \centering
    \begin{tikzpicture}[
        scale=0.8,
        >=Stealth,
        box/.style={
            rounded corners=5pt,
            minimum width=3cm,
            minimum height=0.8cm,
            align=center,
            font=\sffamily\bfseries,
            text=white
        },
        green/.style={box, fill=green!60!black},
        red/.style={box, fill=red},
        arrow/.style={->, thick, draw=blue!60},
        label/.style={font=\sffamily\scriptsize, fill=white, inner sep=2pt}
    ]
    
    % --- Node Placement on a Grid ---
    \node[green] (owners)          at (0, 0)      {Owners};
    \node[green] (availability)    at (5, 0)      {Availability\\usefulness / \ldots};
    \node[green] (countermeasures) at (0, -2.5)   {Countermeasures};
    \node[red]   (vulnerabilities) at (0, -5)     {Vulnerabilities};
    \node[red]   (threats)         at (-4.5, -5)  {Threats};
    \node[red]   (attackers)       at (-9, -5)    {Attackers};
    \node[red]   (risks)           at (4.5, -5)   {Risks};
    \node[green] (assets)          at (4.5, -8)   {Assets};
    
    % --- Edges and Labels ---
    % Owners
    \draw[arrow] (owners) -- node[label, above] {want to maximize} (availability);
    \draw[arrow] (owners) -- node[label, left] {impose} (countermeasures);
    \draw[arrow] (owners.west) to[bend right=45] node[label, left] {want to minimize} (vulnerabilities.west);
    
    % Countermeasures & Vulnerabilities
    \draw[arrow] (countermeasures.250) -- node[label, left] {reduce} (vulnerabilities.110);
    \draw[arrow] (vulnerabilities.70) -- node[label, right] {require} (countermeasures.290);
    
    % Vulnerabilities & Risks
    \draw[arrow] (vulnerabilities.20) to[bend left=15] node[label, above] {lead to} (risks.160);
    \draw[arrow] (vulnerabilities.340) to[bend right=15] node[label, below] {may have} (risks.200);
    
    % Attackers, Threats, Vulnerabilities
    \draw[arrow] (attackers) -- node[label, above] {give rise to} (threats);
    \draw[arrow] (threats) -- node[label, above] {exploit} (vulnerabilities);
    
    % Threats to Risks (Sweeps elegantly under vulnerabilities)
    \draw[arrow] (threats.south) to[bend right=25] node[label, below] {increase} (risks.south);
    
    % Attackers to Assets (L-shape routing under the entire graph)
    \draw[arrow] (attackers.south) |- node[label, near start, right] {want to abuse} (assets.west);
    
    % Risks to Assets
    \draw[arrow] (risks) -- node[label, right] {to} (assets);
    
    % Availability to Assets & Risks
    \draw[arrow] (availability.south) to[bend left=15] node[label, right] {increase} (assets.east);
    \draw[arrow] (availability.250) to[bend right=15] node[label, right] {of} (risks.45);
    
    \end{tikzpicture}
    \caption{Security meta-model definition}
    \label{fig:sec-meta-model}
\end{figure}

\paragraph{How to realize security objectives?}
To minimize risks, there are 3 pillars:
\begin{description}
    \item[Methods] Proactive, Detective, Reactive.
    \subitem[Proactive] Prevent possible issues and flaws at the very beginning.
    \subitem[Detective] Predict/catch issues.
    \subitem[Reactive] Recover assets, repair damage, prosecute offenders.
    \item[Know The Attacker] Types/typical attacks and commonly exploited security issues.
    \item[Existing issues] The attacker does not create security issues and vulnerabilities; they exploit existing ones. 
\end{description}

\paragraph{Security through obscurity} 
\gls{sto} means hiding design and implementation details to gain security. This \textbf{DOES NOT WORK} because it is difficult to keep secrets in coding; if the entire security mechanism relies on a single secret, the whole system is compromised once it is revealed. Also, secret algorithms and unreviewed code mean flaws will not be spotted nor fixed. Security \textbf{and} obscurity can work together, but not ONLY obscurity.

\paragraph{Cryptography} 
Usually, an association between security and cryptography is made; however, many security problems cannot be solved by cryptography. Cryptography and its use have intrinsic limitations. Furthermore, 85\% of \gls{cert} security advisories could not have been solved or prevented with cryptography alone.

\paragraph{When to Start}
\begin{center}
    IMMEDIATELY
\end{center}
Not on the next version.

\chapter{System Development Life Cycle [14/09/2026]}
The \gls{ssdlc} is defined similarly to the traditional \gls{sdlc}, but it focuses on the development of software and applications in a manner that significantly reduces security risks, eliminates security vulnerabilities, and reduces costs. The \gls{ssdlc}, like the \gls{sdlc}, is divided into phases. 

\section{Defining}
In this phase, the product's requirements must be defined and documented, then approved by the customer (or via market analysis) to form a \gls{srs}. 

It produces the working plan that outlines the expected output of the developing team with security milestones, certifications, risk assessments, key security resources, and required third-party resources.

Before outlining the security deliverables of this phase, we must formally define risk and its assessment:

\begin{definition}{Risk}{risk}
    The potential that a given \textbf{threat will exploit vulnerabilities} of an asset or group of assets, thereby having a negative impact on those identified assets. An asset is relevant information, data, a service, or a product.
\end{definition}

\begin{definition}{Risk Assessment}{risk-assessment}
    The process of \textbf{identifying and determining risks} that affect an organization's operations, assets, individuals, other organizations, and the nation.
\end{definition}

\noindent In terms of security, the Defining phase includes:
\begin{enumerate}
    \item \textbf{Gap Analysis}: Defining the security requirements for developers, designers, architects, and \gls{qa} by focusing on secure design principles, security issues, web security, and encryption. This also includes organizing security awareness sessions—both technical-specific training for developers and general sessions for the broader team.
    \item \textbf{Risk Assessment}: The practical execution of the processes defined in \cref{def:risk} and \cref{def:risk-assessment}.
    \item \textbf{Third-Party Software Tracking}: (This will be discussed in later chapters).
\end{enumerate}

\subsection{Threat Modeling}

\begin{definition}{Asset}{asset}
    An asset is relevant data, a service, or a component that is the objective of protection.
\end{definition}

\begin{definition}{Threat}{threat}
    A potential cause of an unwanted incident, which may result in harm to a system or organization (e.g., malicious actors, malware, or natural disasters).
\end{definition}

\begin{definition}{Impact}{impact}
    The magnitude of harm or loss that can be expected if a specific vulnerability is successfully exploited by a threat.
\end{definition}

\begin{definition}{Probability (Likelihood)}{probability}
    The chance or frequency with which a specific threat is expected to successfully exploit a vulnerability.
\end{definition}

\begin{definition}{Quantitative Risk}{quantitative-risk}
    Risk can be calculated as the product of the impact (\cref{def:impact}) resulting from a successful attack on a certain asset (\cref{def:asset}) and the probability (\cref{def:probability}) of that attack occurring.
    \[
        \text{Risk} = \text{Impact} \cdot \text{Probability}
    \]
    The output is a numerical value that is typically mapped to a qualitative severity level (e.g., High, Medium, Low).
\end{definition}

\noindent The threat modeling process is composed of the following steps:
\begin{enumerate}
    \item Identify and characterize the system.
    \item Identify assets (\cref{def:asset}).
    \item Identify threats (\cref{def:threat}).
    \item Identify vulnerabilities (\cref{def:vulnerability}).
    \item Estimate the impact on assets (\cref{def:impact}).
    \item Estimate the probability of exploitation (\cref{def:probability}).
    \item Rank the threats by their calculated risk (\cref{def:quantitative-risk}).
    \item Decide which threats to respond to, using established risk thresholds.
    \item Decide on mitigation and control actions for the selected threats. 
\end{enumerate}
\lecturestop{14/09/2026}

\subsection{Third-Party Software Tracking}
Modern software development heavily relies on code reuse, integrating both open-source and commercial third-party components. Consequently, maintaining a systematic inventory of all third-party software is critical. 

\begin{concept}{Third-Party Software Tracking}{third-party-tracking}
    The systematic process of inventorying, auditing, and monitoring open-source and commercial third-party components integrated into modern software to quickly detect and mitigate upstream vulnerabilities.
\end{concept}

\noindent This tracking allows organizations to audit dependencies and identify when a component becomes vulnerable.

\section{Designing}
During the design phase, architects establish the structural blueprint of the product by creating a formal design specification.

\begin{definition}{Detailed Design Specification (DDS)}{dds}
    A \glsxtrfull{dds} is a comprehensive architectural blueprint that defines system structure and behavior based on several key factors: risk assessments (\cref{def:risk-assessment}), product robustness, design modularity, budget limitations, and time constraints.
\end{definition}

\noindent The resulting architecture is partitioned into the overall project structure and the specific architecture of individual modules. Furthermore, it is crucial to precisely define the data flows between the system, external entities, and third-party modules.

\begin{concept}{Secure Architectural Principles}{secure-arch-principles}
    The fundamental rules for secure code development include:
    \begin{itemize}
        \item \textbf{Modularity:} Encapsulating functionality into distinct components, modules, and functions.
        \item \textbf{Isolation:} Ensuring each component functions independently to contain potential failures.
        \item \textbf{Reusability:} Combining modularity and isolation safely.
        \item \textbf{Hierarchical Organization:} Structuring system dependencies logically.
        \item \textbf{Simplicity:} Keeping designs clean and understandable to minimize hidden flaws.
    \end{itemize}
\end{concept}

\paragraph{Design Review} The design requires rigorous review from a developer's perspective. This process must incorporate a full evaluation of security aspects alongside functional features. During this security review, developers must proactively adopt an attacker's mindset to discover potential feature vulnerabilities. If a flaw is identified, developers should rely on established secure design checklists that enforce controls such as the mandatory logging of security events and failures.

\section{Implementation}
Once the \gls{dds} is finalized, the implementation phase begins, translating specifications into code. If the design phase has been conducted rigorously, coding proceeds with structured clarity. The choice of programming language depends heavily on the specific domain and nature of the software being developed.

\noindent From a security perspective, we recall that a bug represents a deviation from design intent (\cref{def:bug}), while a vulnerability is a causal chain starting with a bug and resulting in an error that, if exploited, leads to a security failure (\cref{def:vulnerability}). Many security incidents stem from vulnerabilities introduced by bugs, defects, or design flaws. To mitigate these risks during coding, developers should:
\begin{itemize}
    \item Follow established security guidelines;
    \item Critically evaluate the security implications of implementation choices;
    \item Reuse trusted and vetted code;
    \item Write high-quality, readable, and maintainable code;
    \item Conduct thorough code reviews;
    \item Verify and test both individual components and the overall system.
\end{itemize}

\begin{concept}{Secure Coding Guidelines}{secure-coding-guidelines}
    General secure coding guidelines encompass core industry best practices, such as:
    \begin{itemize}
        \item Treating all input data as untrusted;
        \item Enforcing strict allowlisting (accepting only valid characters and formats);
        \item Strictly separating data from code (e.g., to prevent injection attacks);
        \item Protecting sensitive information through proper encryption and secure storage;
        \item Never trusting the client-side state.
    \end{itemize}
    Frameworks can significantly enhance the adoption of these practices, as framework code is actively maintained and audited for vulnerabilities.
\end{concept}

\paragraph{Code Review} Numerous security defects can be detected directly in source code. 

\begin{definition}{Code Review}{code-review}
    A manual or automated code inspection process focused on identifying implementation-level bugs, anti-patterns, and common code security pitfalls by leveraging standardized review checklists.
\end{definition}

\paragraph{Overall View} As previously noted, during the implementation phase, primary security verification relies on \textbf{static analysis} and \textbf{vulnerability scanning}, because the complete system is not yet operational or ready to execute in a live environment.

\section{Testing}
Testing must not be treated as a single, isolated phase at the end of development; rather, it should be performed iteratively across all stages of the product development lifecycle. Identified defects must be systematically reported, tracked, fixed, and retested. This feedback loop is repeated continuously until the software quality fully satisfies the requirements specified in the \gls{srs}. 

\noindent From an advanced security perspective, testing incorporates dynamic and active methodologies:

\begin{concept}{Dynamic Application Security Testing (DAST)}{dast}
    Analyzing the running application in real-time to catch runtime anomalies, memory leakage, and execution errors.
\end{concept}

\begin{concept}{Fuzzing}{fuzzing}
    An automated testing technique that feeds semi-random, malformed, or unexpected inputs into the system to expose hidden crash bugs and memory corruption vulnerabilities.
\end{concept}

\begin{concept}{Penetration Testing (Pentest)}{pentest}
    Engaging third-party security experts to actively simulate real-world attacks against the system to discover exploitable weaknesses.
\end{concept}

\section{Deployment and Maintenance}
Once the software passes rigorous verification and testing phases, it transitions into deployment and long-term maintenance. 

\begin{itemize}
    \item \textbf{Alpha/Beta Releases:} Initial restricted releases (alpha/beta) deployed to collect real-world user feedback on software usage, usability, and unanticipated edge cases.
    \item \textbf{Final Security Milestones:} Concluding the deployment pipeline with a final gap analysis and comprehensive security audits before pushing updates to production.
\end{itemize}

\subsection{The Economics of Bug Fixing}
A foundational principle in software engineering and security is that \textbf{bugs (including security vulnerabilities) should be fixed as soon as possible}. 

\begin{observation}{Cost of Fixing Defects}{cost-defects}
    According to studies by the IBM Systems Science Institute, the relative cost of fixing a software defect increases exponentially the later it is discovered in the life cycle. Fixing a vulnerability during the requirements or design phase is vastly cheaper than fixing it after deployment in a production environment. 
\end{observation}

\noindent Consequently, security cannot be bolted on as an afterthought once implementation is complete; doing so is fundamentally ineffective and high-risk. Instead, security must be integrated natively into every phase of the \gls{ssdlc}.

\subsection{DevOps and DevSecOps}
Modern software engineering heavily relies on \textbf{DevOps} (a compound of ``Development'' and ``Operations''), which aims to bridge the gap between building software and running it reliably in production.

\begin{concept}{DevOps Practices}{devops-practices}
    DevOps is structured around iterative delivery cycles (such as Scrum and Agile) designed for rapid feature delivery and quick feedback loops (promoting a ``fail fast, fix fast'' culture). Its core practices include:
    \begin{itemize}
        \item \textbf{Continuous Integration (CI):} Developers regularly merge code changes into a central repository multiple times a day, triggering automated builds and test suites to catch bugs immediately.
        \item \textbf{Continuous Delivery (CD):} Code changes that pass automated testing are automatically prepared and staged for immediate production deployment, ensuring the repository always maintains a deployment-ready artifact.
        \item \textbf{Supporting Practices:} Utilization of microservices, infrastructure as code (IaC), centralized monitoring/logging, and cross-functional team communication.
    \end{itemize}
\end{concept}

\noindent When security practices, automated security testing, and vulnerability checks are fully embedded into this DevOps pipeline, the methodology evolves into \textbf{DevSecOps}—ensuring that rapid velocity does not compromise safety, integrity, or compliance.

\chapter{\glsfmtfull{cwe}}
Recalling our earlier definitions of a bug (\cref{def:bug}) and a weakness (\cref{def:vulnerability}), we observe that software weaknesses are fundamentally rooted in code implementation choices. However, recognizing typical code patterns that lead to security issues is insufficient on its own; vulnerability manifestation depends heavily on the execution environment, compiler optimizations, and system context. Consequently, mastery of secure programming guidelines and knowing precisely what patterns to target during security analysis and testing are critical.

\noindent As previously introduced, the \gls{cwe} provides a standardized, publicly available catalog of the most common software and hardware weaknesses. Each \gls{cwe} entry is characterized by several structured fields, including a unique identifier, abstraction level, and a hierarchical organization (e.g., *parent of*, *child of*, *member of*). Furthermore, lists like the \textit{CWE Top 25} highlight the most dangerous software weaknesses across various systems by scoring entries based on how frequently they map to known \glspl{cve} and averaging their real-world severity.

\section{CWE-416:\glsfmtfull{uaf}}
A \textbf{\gls{uaf}} vulnerability occurs when a product reuses or references memory after it has been freed. Subsequent operations can result in that memory region being reallocated and assigned to a different pointer. Any operations performed via the original (dangling) pointer are no longer valid, as the memory now belongs to data processed by completely different components.

\paragraph{Prevention} To prevent \gls{uaf} vulnerabilities, developers should immediately set pointers to \texttt{NULL} after freeing them, or utilize programming languages equipped with automatic memory management (garbage collection).

\paragraph{Detection} Fuzzing, automated \gls{sast}, and dynamic runtime analysis.

\paragraph{Related Attacks} Code injection (\gls{capec}-242), Buffer Overflow (\gls{capec}-100).

\paragraph{Notes} \gls{uaf} is not restricted strictly to memory pointers; it can also affect abstract resources such as file handlers, sockets, and network connections. Hence, \gls{uaf} vulnerabilities can occasionally manifest even in memory-managed languages if underlying resource lifecycles are mismanaged.

\paragraph{Real-World Examples (\glspl{cve})}
\begin{itemize}
    \item \textbf{CVE-2022-20141:} Insufficient resource locking in the Android system kernel (CWE-413) leading to a \gls{uaf} condition;
    \item \textbf{CVE-2022-2621:} Two concurrent threads in a web browser attempt to share a resource (CWE-366), allowing one thread to destroy the resource before the other completes execution;
    \item \textbf{CVE-2020-6819:} A race condition (CWE-362) leading to a \gls{uaf} exploit vector;
    \item \textbf{CVE-2010-4168:} A \gls{uaf} triggered by closing a network connection while data transmission is still actively ongoing.
\end{itemize}

\section{CWE-122: Heap-based Buffer Overflow}
A heap overflow condition is a buffer overflow where the buffer that can be overwritten is allocated in the \underline{heap portion of memory}, generally meaning that the buffer was allocated dynamically.

\paragraph{Prevention} Using automatic buffer overflow detection mechanisms, implementing strict bounds checking on input, and avoiding the use of dangerous functions such as \texttt{gets} (ensuring safe equivalents are used instead).

\paragraph{Detection} Fuzzing, automated \gls{sast}.

\paragraph{Related Attacks} Forced Integer Overflow (\gls{capec}-92), Overflow Buffers (\gls{capec}-100).

\paragraph{Real-World Examples (\glspl{cve})}
\begin{itemize}
    \item \textbf{CVE-2021-29529:} A machine-learning product suffers from a heap-based buffer overflow (CWE-122) when integer-oriented bounds are calculated using \texttt{ceiling()} and \texttt{floor()} on floating-point values (CWE-1339);
    \item \textbf{CVE-2021-43537}: In a web browser, an unsigned 64-bit integer is forcibly cast to a 32-bit integer (CWE-681), potentially leading to an integer overflow (CWE-190) and subsequent heap memory corruption (CWE-122);
    \item \textbf{CVE-2007-4268}: An integer signedness error (CWE-195) bypasses signed comparison, leading to a heap overflow;
    \item \textbf{CVE-2009-2523}: A product fails to handle non-\texttt{NULL}-terminated input strings (CWE-170), leading to buffer over-read or heap-based buffer overflow.
\end{itemize}

\section{CWE-787: Out-of-bounds Write}
The product writes data past the intended limit of a buffer, often resulting in memory corruption or arbitrary code execution outside the buffer's valid boundaries.

\paragraph{Prevention} Implement rigorous checks on size and initial position before utilizing the buffer; replace unbounded copy functions with safer alternatives that support length arguments (e.g., replacing \texttt{strcpy} with \texttt{strncpy}), and implement custom bounds checks if safe APIs are unavailable.

\paragraph{Detection} Automated \gls{sast}, automated dynamic analysis with fuzzing and fault-injection-based testing.

\paragraph{Notes} The classic ``buffer overflow'' is a specific type of out-of-bounds write that occurs within structured buffers (e.g., arrays).

\paragraph{Related Attacks} Overflow Buffers (\gls{capec}-100), Code Injection (\gls{capec}-242).

\paragraph{Related \glspl{cwe}} \textbf{CWE-119}: Improper Restriction of Operations within the Bounds of a Memory Buffer (performing read or write operations outside a buffer's intended boundaries).

\paragraph{Real-World Examples (\glspl{cve})}
\begin{itemize}
    \item \textbf{CVE-2021-21220}: Insufficient input validation (CWE-20) in a browser component allows heap corruption;
    \item \textbf{CVE-2020-17087}: Integer truncation (CWE-197) causes a small buffer allocation, leading to an out-of-bounds write;
    \item \textbf{CVE-2019-1010006}: Compiler optimization (CWE-733) improperly removes or modifies code designed to detect integer overflow;
    \item \textbf{CVE-2020-0968}: A remote code execution vulnerability in the Internet Explorer scripting engine handling objects in memory (Scripting Engine Memory Corruption Vulnerability).
\end{itemize}

\section{CWE-78: OS Command Injection}
The product constructs an OS command using externally-influenced input, but it does not (or incorrectly) neutralize special elements that could modify the intended OS command.

\paragraph{Prevention} Use dedicated, built-in language libraries or APIs rather than spawning external processes to recreate the desired functionality. Furthermore, security checks performed on the client side \textbf{MUST} be strictly replicated on the server side. When executing commands, consider passing user-supplied arguments via standard input (stdin) or securely parameterized APIs rather than appending them directly to command-line execution strings.

\paragraph{Detection} Automated static analysis is often insufficient on its own. A robust dynamic analysis approach—incorporating fuzzing and fault-injection-based testing—is much more adequate for discovering execution context vulnerabilities.

\paragraph{Related Attacks} OS Command Injection (CAPEC-88), Command Line Execution through SQL Injection (CAPEC-108).

\paragraph{Real-World Examples (\glspl{cve})}
\begin{itemize}
    \item \textbf{CVE-2020-10987}: OS command injection in a Wi-Fi router; allows remote attackers to execute arbitrary system commands via the \texttt{deviceName} POST parameter.
    \item \textbf{CVE-2020-9054}: Improper input validation (CWE-20) in a username parameter, ultimately leading to OS command injection.
    \item \textbf{CVE-1999-0067}: A CGI program fails to neutralize the ``\texttt{|}'' metacharacter when invoking a phonebook application.
    \item \textbf{CVE-2003-0041}: An FTP client fails to filter ``\texttt{|}'' from filenames returned by the server, allowing for client-side OS command injection.
\end{itemize}

\section{CWE-20: Improper Input Validation}
The product receives input values but does not (or incorrectly) validates them. Input validation is a fundamental and frequently used technique for checking potentially dangerous inputs to ensure they are safe for processing. In the absence of proper validation, unintended input values can maliciously alter the program's control flow.

\paragraph{Prevention} Implement a robust input validation framework. Adopt a zero-trust approach: assume all inputs are malicious, check all potentially untrusted inputs against a strict allowlist (positive validation), and enforce strict type, length, and format constraints.

\paragraph{Detection} Automated and manual static analysis combined with dynamic analysis (such as fuzzing).

\paragraph{Related Attacks} Command Line Execution through SQL Injection (CAPEC-108), Format String Injection (CAPEC-135), Input Data Manipulation (CAPEC-153), Blind SQL Injection (CAPEC-7), Server-Side Request Forgery (SSRF) (CAPEC-664), Cross-Site Scripting (XSS) (CAPEC-63).

\paragraph{Real-World Examples (\glspl{cve})}
\begin{itemize}
    \item \textbf{CVE-2024-37032}: Large language model (LLM) management tool does not validate the format of a digest value from a private, untrusted model registry.
    \item \textbf{CVE-2020-9054}: Improper input validation in username parameter, leading to OS command injection.
    \item \textbf{CVE-2021-30860}: Improper input validation leads to integer overflow in mobile OS, as exploited in the wild.
    \item \textbf{CVE-2021-22205}: Backslash followed by a newline can bypass a validation step, leading to eval injection, as exploited in the wild.
\end{itemize}

\section{CWE-502: Deserialization of Untrusted Data}
In numerous occasions, it is convenient to serialize objects for communication or storage. This vulnerability occurs when the product deserializes untrusted data without sufficiently verifying that the resulting data will be valid and safe. Serialized data (or code) can be maliciously altered to bypass standard accessor functions if there is no cryptographic mechanism protecting its integrity.

\paragraph{Prevention} When processing serialized data, extract the primitive fields to populate a new, safe object rather than directly instantiating complex types from the untrusted stream. Furthermore, ensure the integrity of the data has not been compromised by enforcing digital signatures or \glspl{hmac} (Hash-based Message Authentication Codes) before attempting to deserialize it. 

\paragraph{Detection} Automated static analysis is generally sufficient to detect the presence of unchecked deserialization calls (e.g., insecure uses of \texttt{ObjectInputStream} in Java or \texttt{pickle} in Python).

\paragraph{Related Attacks} Object Injection (CAPEC-586).

\paragraph{Real-World Examples (\glspl{cve})}
\begin{itemize}
    \item \textbf{CVE-2019-12799}: Bypass of untrusted deserialization issue (CWE-502) by using an assumed-trusted class.
    \item \textbf{CVE-2015-8103}: Deserialization issue in commonly-used Java library allows remote execution.
    \item \textbf{CVE-2013-1465}: Use of PHP unserialize function on untrusted input allows attacker to modify application configuration.
\end{itemize}

\section{CWE-22: Path Traversal}
The full name is: ``Improper Limitation of a Pathname to a Restricted Directory.'' This means the product uses external input to construct a pathname for a file or a directory located within a restricted parent directory. This is a classic situation of improper neutralization of input, since the inserted path can resolve to a different location than the one intended (e.g., by using traversal separators such as \texttt{../} or \texttt{/}).

\paragraph{Prevention} Assume all input is malicious. Store library, include, and utility files outside the web document root whenever possible. When the set of acceptable objects (filenames, URLs, etc.) is known and limited, create a strict mapping from fixed, safe input values to the actual target filenames.

\paragraph{Detection} Automated and manual static analysis, coupled with automated dynamic analysis and vulnerability scanning.

\paragraph{Real-World Examples (\glspl{cve})}
\begin{itemize}
    \item \textbf{CVE-2024-4315}: Local File Inclusion (LFI) in \texttt{lollms} due to insufficient path sanitization of Windows-style backward slashes.
    \item \textbf{CVE-2019-20916}: Path traversal in Python \texttt{pip} allowing arbitrary file overwrite via a crafted wheel (\texttt{.whl}) filename.
    \item \textbf{CVE-2022-31503}: Absolute path traversal in the \texttt{orchest} GitHub repository due to unsafe use of the Flask \texttt{send\_file} function.
\end{itemize}

\section{CWE-918: Server-Side Request Forgery (SSRF)}
Assuming there exists a server-to-server or service-to-service trust relationship, a security vulnerability in a server could allow an attacker to cause the server-side application to make forged requests to an unintended location. The attacker can bypass access control measures by means of the trusted server, using it as a proxy to conduct activities such as port scanning on hosts within an internal network.

\noindent Note that this attack is not limited to targeting other servers; it can also be used to access restricted local files residing on the original server itself.

\paragraph{Prevention} Limit outbound traffic from the original application server. Build a strict allowlist (not a denylist) of trusted domains and IP addresses, and rigorously sanitize user input to prevent potential execution risks.

\paragraph{Detection} Automated static analysis.

\paragraph{Related Attacks} Server-Side Request Forgery (CAPEC-664).

\paragraph{Real-World Examples (\glspl{cve})}
\begin{itemize}
    \item \textbf{CVE-2021-26855}: Microsoft Exchange Server SSRF vulnerability (ProxyLogon) allowing unauthenticated attackers to send arbitrary HTTP requests.
    \item \textbf{CVE-2021-21973}: SSRF in VMware vSphere Client due to improper validation of URLs in a vCenter Server plugin, leading to information disclosure.
    \item \textbf{CVE-2016-4029}: WordPress SSRF bypass achieved by using octal and hexadecimal IP address formats to evade intranet address checks.
    \item \textbf{CVE-2002-1484}: DB4Web server verbose debugging feature allows remote attackers to use the server as a proxy to conduct TCP port scans.
\end{itemize}

\section{Other \glsfmtshortpl{cwe} in the Top 25}
\begin{itemize}
    \item \textbf{CWE-79}: Improper Neutralization of Input During Web Page Generation (Cross-Site Scripting / XSS).
    \item \textbf{CWE-89}: Improper Neutralization of Special Elements used in an SQL Command (SQL Injection).
    \item \textbf{CWE-352}: Cross-Site Request Forgery (CSRF).
    \item \textbf{CWE-476}: NULL Pointer Dereference.
    \item \textbf{CWE-190}: Integer Overflow or Wraparound.
    \item \textbf{CWE-362}: Concurrent Execution using Shared Resource with Improper Synchronization (Race Condition).
    \item \textbf{CWE-306}: Missing Authentication for Critical Function.
    \item \textbf{CWE-94}: Improper Control of Generation of Code (Code Injection).
\end{itemize}

\chapter{The \glsfmtshort{owasp} Top 10}
As established in previous sections, a single \gls{cwe} or a chain of \glspl{cwe} can lead to a specific \gls{cve}. While each represents a distinct security risk, analyzing them individually is inefficient at scale. Fortunately, we can group them into broader ``security issue types.'' 

Characterization and classification are the foundational steps in vulnerability investigation. Grouping issues allows security teams to focus on the most relevant and impactful threats, design more effective testing strategies, improve the \gls{sdlc}, and drive prioritized bug-fixing. Before diving into the \gls{owasp} Top 10 itself, it is important to understand the broader context of how and why these classifications are made.

\section{Classification Criteria for Security Errors}
To properly classify software security errors, several theoretical and practical criteria are used by the industry and academia. Classification goes beyond simple impact; it analyzes the root cause, location, and timing of the defect.

\begin{description}
    \item[Severity] The technical impact that the bug or error has on the system. It answers the questions: How severe is the vulnerability? How critical is the potential damage? (Typically measured on a scale such as High, Medium, Low, or via the \gls{cvss}).
    \item[Priority] How quickly the issue should be resolved based on business context and exploitability. A vulnerability might have high severity but low priority if it exists in an unreachable, isolated internal network.
    \item[Nature and Domain] The technical characteristics of the bugs. Are they functional defects, pure security flaws, or architectural design errors? Furthermore, are they code-specific (e.g., C++ memory leaks) or technology-related (e.g., SQL misconfigurations)?
    \item[Genesis (Intent)] Security errors are often classified by their genesis: whether they are \textit{unintentional} (e.g., logic bugs, typos, memory management errors by honest developers) or \textit{intentional} (e.g., malicious logic, backdoors, or logic bombs inserted by a rogue insider).
    \item[Time of Introduction] Flaws can be categorized by the exact phase of the \gls{sdlc} in which they were introduced—such as requirement gathering, architectural design, implementation, or deployment/maintenance. Fixing an error introduced during the design phase is notoriously more expensive than fixing an implementation bug.
\end{description}

\section{Context: Prominent Industry Taxonomies}
Several standardized taxonomies exist to help developers and testers navigate the vast landscape of security errors. While this section focuses heavily on the \gls{owasp} Top 10, it is part of a broader ecosystem of security classifications:

\paragraph{The \gls{cwe} Top 25} 
Maintained by MITRE, the \gls{cwe} Top 25 list highlights the most dangerous software weaknesses across all applications and systems, regardless of the specific domain. These weaknesses are scored based on two main metrics: \textit{prevalence} (how often a \gls{cwe} is mapped to a real-world \gls{cve}) and \textit{severity} (the potential for damage). A distinct subset, the \gls{cwe} Top 10, is specifically aggregated based on actively \textbf{known exploited vulnerabilities} in the wild.

\paragraph{The 24 Deadly Sins of Software Security} 
Coined by security experts Michael Howard, David LeBlanc, and John Viega, the ``24 Deadly Sins'' is a highly practical, developer-centric taxonomy. Rather than focusing purely on abstract theoretical models, it groups vulnerabilities by programming flaws that developers make every day. The sins are categorized into four distinct areas: \textit{Web Application Flaws} (e.g., \gls{xss}, \gls{sqli}), \textit{Implementation Flaws} (e.g., Buffer overflows in C/C++, integer overflows), \textit{Cryptographic Flaws} (e.g., using weak algorithms or hardcoded keys), and \textit{Networking/Operational Flaws}.

\paragraph{The \gls{owasp} Top 10} 
The \gls{owasp} Top 10 is the globally recognized standard awareness document for developers and web application security. Unlike the \gls{cwe} list, which covers all software (including \glspl{os} and embedded systems), \gls{owasp} focuses strictly on web application vulnerabilities (e.g., Broken Access Control, Cryptographic Failures, Injection). It represents a broad consensus about the most critical security risks to web applications.

It is important to note that the \gls{owasp} Top 10 is updated approximately every four years (the latest edition being 2025). Between iterations, vulnerabilities are sometimes merged, split, or re-ranked based on current threat intelligence, and entirely new typologies can be introduced.

\section{A01-2025: Broken Access Control}
Access control normally enforces policies such that users cannot act outside of their intended permissions. Abnormalities in this mechanism lead to unauthorized information disclosure, modification, or destruction of data by allowing the execution of actions outside the user's granted limits.

\paragraph{Prevention} A \textit{deny-by-default} policy should be implemented. Furthermore, access control is only effective in trusted server-side code and/or server-less \glspl{api}, where the attacker cannot modify the access control checks or metadata. Rate limiting on the \gls{api} controllers must be enforced to minimize the danger of automated attacks. Finally, stateful session identifiers must be invalidated on the server immediately after logout.

\paragraph{Real-World Examples (\glspl{cwe})}
\begin{itemize}
    \item \textbf{CWE-22}: Improper Limitation of a Pathname to a Restricted Directory (Path Traversal)
    \item \textbf{CWE-23}: Relative Path Traversal
    \item \textbf{CWE-352}: \glsxtrfull{csrf}
    \item \textbf{CWE-601}: URL Redirection to Untrusted Site (Open Redirect)
    \item \textbf{CWE-566}: Authorization Bypass Through User-Controlled SQL Primary Key
    \item \textbf{CWE-377}: Insecure Temporary File
    \item \textbf{CWE-200}: Exposure of Sensitive Information to an Unauthorized Actor
    \item \textbf{CWE-918}: \glsxtrfull{ssrf}
\end{itemize}

\noindent \textit{Notice that diverse \glspl{cwe}—ranging from path traversal to server-side request forgery—can all culminate in the same overarching security risk of broken access control.}

\section{A02-2025: Security Misconfiguration}
This category refers to all situations where there is missing (or improperly configured) security hardening across any part of the application stack. It includes scenarios where unnecessary features are enabled or installed, default accounts and passwords are still active and unchanged, error handling reveals sensitive stack traces, or the software itself is out of date.

\paragraph{Prevention} To prevent this, a secure, repeatable installation process should be implemented. Additionally, routine automated tasks to review, audit, and upgrade these installations and configurations must be scheduled regularly.

\paragraph{Real-World Examples (\glspl{cwe})}
\begin{itemize}
    \item \textbf{CWE-13}: ASP.NET Misconfiguration: Password in Configuration File
    \item \textbf{CWE-15}: External Control of System or Configuration Setting
    \item \textbf{CWE-537}: Java Runtime Error Message Containing Sensitive Information
    \item \textbf{CWE-756}: Missing Custom Error Page
\end{itemize}

\section{A03-2025: Software Supply Chain Failures}
This category groups all failures that cause breakdowns in the process of building, distributing, or updating software. These breakdowns are often caused by vulnerabilities or malicious changes in third-party code, build tools, or other external dependencies.

\paragraph{Prevention} To prevent this, organizations must maintain and manage a comprehensive \glsxtrfull{sbom} and continuously monitor libraries and third-party components. Developers must track all direct and transitive dependencies, actively removing unused or unnecessary features, components, and files. Finally, there must be constant monitoring of \glspl{cve} via the \gls{nvd} for newly discovered vulnerabilities.

\paragraph{Real-World Examples (\glspl{cwe})}
\begin{itemize}
    \item \textbf{CWE-1104}: Use of Unmaintained Third-Party Components
    \item \textbf{CWE-1395}: Dependency on Vulnerable Third-Party Component
    \item \textbf{CWE-1329}: Reliance on Component That is Not Updateable
\end{itemize}

\section{A04-2025: Cryptographic Failures}
This category encompasses all violations of the protection needs for sensitive data, such as passwords, credit card numbers, health records, personal information, and business secrets. All of the aforementioned data require extra protection, especially since they are usually governed by strict privacy laws and regulations like the \glsxtrfull{gdpr} or the \glsxtrfull{pcidss}.

\paragraph{Prevention} Ensure the proper encryption of all sensitive data both \textit{at rest} and \textit{in transit}. Furthermore, ensure the correct and secure configuration of all applied cryptographic mechanisms (e.g., using strong, modern algorithms, proper key management, and avoiding deprecated hash functions).

\paragraph{Real-World Examples (\glspl{cwe})}
\begin{itemize}
    \item \textbf{CWE-261}: Weak Encoding for Password
    \item \textbf{CWE-319}: Cleartext Transmission of Sensitive Information
    \item \textbf{CWE-321}: Use of Hard-coded Cryptographic Key
\end{itemize}

\section{A05-2025: Injection}
An application is vulnerable to injection attacks in any of the following cases: user data is neither validated nor sanitized; dynamic queries are constructed directly without parameterized calls or context-aware escaping; hostile data is used within \glsxtrfull{orm} search parameters to extract additional sensitive records; or hostile data is directly used or concatenated into interpreters.

\paragraph{Prevention} To prevent an injection, untrusted data must be kept strictly separate from commands and queries. Robust server-side input validation must be implemented (supplemented by client-side validation for improved \glsxtrfull{ux}). Furthermore, all special characters must be properly escaped, and the permissions or returned records of \gls{sql} queries should be strictly limited (e.g., using \texttt{LIMIT} and enforcing the principle of least privilege).

\paragraph{Real-World Examples (\glspl{cwe})}
\begin{itemize}
    \item \textbf{CWE-20}: Improper Input Validation
    \item \textbf{CWE-78}: Improper Neutralization of Special Elements used in an \gls{os} Command (OS Command Injection)
    \item \textbf{CWE-79}: Improper Neutralization of Input During Web Page Generation (\glsxtrfull{xss})
    \item \textbf{CWE-89}: Improper Neutralization of Special Elements used in an \gls{sql} Command (\glsxtrfull{sqli})
    \item \textbf{CWE-94}: Improper Control of Generation of Code (Code Injection)
    \item \textbf{CWE-116}: Improper Encoding or Escaping of Output
\end{itemize}

\section{A06-2025: Insecure Design}
The weaknesses in this category are related to ``Missing or Ineffective Control Design.'' Insecure design cannot be fixed by a simple code patch; it requires architectural changes. Secure design encompasses:
\begin{itemize}
    \item \textbf{Requirements and Resource Management}: Research core business requirements and negotiate security requirements concerning the \gls{cia} triad and authenticity for all assets.
    \item \textbf{Secure Design}: Constantly evaluate threats, ensuring code is robustly designed and tested against business logic flaws.
    \item \textbf{Secure Development Lifecycle}: Integrate security controls throughout the entire development process (via a \gls{ssdlc}).
\end{itemize}

\paragraph{Prevention} To prevent vulnerabilities caused by insecure design, organizations must adopt a secure design methodology, conduct threat modeling, and properly integrate security controls at the architectural design stage.

\paragraph{Real-World Examples (\glspl{cwe})}
\begin{itemize}
    \item \textbf{CWE-73}: External Control of File Name or Path
    \item \textbf{CWE-256}: Plaintext Storage of a Password
    \item \textbf{CWE-444}: Inconsistent Interpretation of \gls{http} Requests (\gls{http} Request/Response Smuggling)
    \item \textbf{CWE-451}: User Interface (\gls{ui}) Misrepresentation of Critical Information
    \item \textbf{CWE-501}: Trust Boundary Violation
\end{itemize}

\section{A07-2025: Authentication Failures}
In these vulnerabilities, an attacker tricks the application into recognizing an invalid or untrusted identity as a legitimate user. This category encompasses flaws in identity confirmation, credential verification, and session management—all of which are critical for protecting against authentication-related attacks.

\paragraph{Prevention} To prevent vulnerabilities in this category, robust identification, authentication, and session handling mechanisms must be enforced. Recommended practices include implementing \glsxtrfull{mfa}, preventing brute-force attacks via rate limiting, and blocking the reuse of stolen credentials. Furthermore, strict password complexity requirements must be enforced. Finally, all authentication checks and session tokens must be generated and validated strictly within a secure server-side backend.

\paragraph{Real-World Examples (\glspl{cwe})}
\begin{itemize}
    \item \textbf{CWE-259}: Use of Hard-coded Password
    \item \textbf{CWE-306}: Missing Authentication for Critical Function
    \item \textbf{CWE-384}: Session Fixation
    \item \textbf{CWE-1391}: Use of Weak Credentials
    \item \textbf{CWE-613}: Insufficient Session Expiration
\end{itemize}
\backmatter
\chapter{Appendix}

% Generate automated lists of definitions, concepts, and theorems
\tcblistof[\section]{deflist}{List of Definitions}
\tcblistof[\section]{cptlist}{List of Concepts}
\tcblistof[\section]{thmlist}{List of Theorems}

\clearpage
% Set a nice gap between the two columns
\setlength{\columnsep}{1cm}

% Span the section title across both columns
\begin{multicols}{2}[\section{List of Acronyms}]
    % \raggedright
    % Disable the default glossary title so it doesn't print twice
    \renewcommand{\glossarysection}[2][]{}
    \printglossary[type=\acronymtype]
\end{multicols}

\end{document}